\chapter{Balance Sheet, Income Statement \& Fraud Vulnerabilities}

% ----------------------------------------------------------------------
% PART 1: EXECUTIVE ROADMAP & LEARNING OBJECTIVES
% ----------------------------------------------------------------------
\section{Executive Roadmap \& Learning Objectives}

The balance sheet and income statement form the twin structural pillars of financial disclosure. While the balance sheet captures the cumulative stock of enterprise assets, obligations, and net worth at a discrete balance point, the income statement reports the dynamic economic flow of revenues and expenses incurred to generate profit over an operating period. For the forensic investigator, these two statements represent the primary canvas upon which corporate earnings manipulation and asset misrepresentation are executed.

\begin{tcolorbox}[colback=primarylight, colframe=primarydark, title={\textbf{Core Learning Competencies}}]
Upon mastering this chapter, students will be equipped to:
\begin{itemize}
    \item \textbf{Dissect Statement Hierarchies:} Formulate the multi-step income statement waterfall (Gross Profit, Operating Profit/EBIT, EBITDA, Net Income) and the classified balance sheet liquidity ordering.
    \item \textbf{Map the ACFE Fraud Architecture:} Classify enterprise misconduct across the three branches of the Association of Certified Fraud Examiners (ACFE) Fraud Tree: Corruption, Asset Misappropriation, and Financial Statement Fraud.
    \item \textbf{Apply the Fraud Triangle and Diamond Models:} Evaluate corporate fraud risk across the four psycho-structural dimensions: Incentive/Pressure, Opportunity, Rationalization, and Executive Capability.
    \item \textbf{Detect Common Statement Manipulation Typologies:} Identify cookie-jar reserve adjustments, aggressive expense capitalization, channel stuffing, and premature revenue recognition.
    \item \textbf{Execute Computational Statement Forensics:} Implement automated vertical (common-size) and horizontal (trend) ratio analysis engines in Python to flag structural financial anomalies.
\end{itemize}
\end{tcolorbox}

\begin{table}[htbp]
\centering
\small
\begin{tabularx}{\textwidth}{llX}
\toprule
\textbf{Financial Construct} & \textbf{Symbol / Variable} & \textbf{Accounting Identity / Operational Equation} \\
\midrule
Gross Profit & $GP$ & $\text{Net Revenue} - \text{COGS}$ \\
Operating Profit & $\text{EBIT}$ & $GP - \text{SG\&A} - \text{R\&D} - \text{Depreciation}_{\text{operating}}$ \\
Operational EBITDA & $\text{EBITDA}$ & $\text{EBIT} + \text{Depreciation} + \text{Amortization}$ \\
Net Operating Working Capital & $NOWC$ & $(\text{Current Assets} - \text{Cash}) - (\text{Current Liabilities} - \text{Short-Term Debt})$ \\
Vertical Metric (Balance Sheet) & $v_{i,t}^{\text{BS}}$ & $\frac{\text{Line Item}_{i,t}}{\text{Total Assets}_t} \times 100\%$ \\
Vertical Metric (Income Stmt) & $v_{i,t}^{\text{IS}}$ & $\frac{\text{Line Item}_{i,t}}{\text{Total Net Revenue}_t} \times 100\%$ \\
Horizontal Metric & $h_{i,t}$ & $\frac{\text{Line Item}_{i,t} - \text{Line Item}_{i,t-1}}{\text{Line Item}_{i,t-1}} \times 100\%$ \\
\bottomrule
\end{tabularx}
\caption{Mathematical Notation Guide for Financial Statement Decomposition.}
\label{tab:w2_notation}
\end{table}

% ----------------------------------------------------------------------
% PART 2: MATHEMATICAL & COMPUTATIONAL FOUNDATIONS
% ----------------------------------------------------------------------
\section{Mathematical \& Computational Foundations}

\subsection{The Classified Balance Sheet Liquidity Model}
In accordance with standard accounting frameworks (IFRS and US GAAP), balance sheet line items are organized in decreasing order of \index{liquidity}\textbf{liquidity} (for assets) and temporal maturity (for liabilities):
\begin{align}
\text{Total Assets}_t &= \underbrace{\text{Cash}_t + \text{MktSec}_t + \text{AR}_t + \text{Inv}_t + \text{Prepaid}_t}_{\text{Current Assets } (CA_t)} + \underbrace{\text{PP\&E}_t + \text{Intangibles}_t + \text{LT-Inv}_t}_{\text{Non-Current Assets } (NCA_t)} \\
\text{Total Claims}_t &= \underbrace{\text{AP}_t + \text{NotesPay}_t + \text{AccruedExp}_t}_{\text{Current Liabilities } (CL_t)} + \underbrace{\text{BondsPay}_t + \text{LT-Debt}_t}_{\text{Non-Current Liabilities } (NCL_t)} + \text{Equity}_t
\end{align}

Working Capital represents the enterprise's short-term liquidity buffer:
\begin{equation}
WC_t = CA_t - CL_t, \quad \text{Current Ratio}_t = \frac{CA_t}{CL_t}, \quad \text{Quick Ratio}_t = \frac{CA_t - \text{Inv}_t}{CL_t}
\end{equation}

\subsection{The Multi-Step Profit Decomposition}
Unlike a single-step statement, the multi-step income statement isolates gross, operational, financing, and tax dimensions:
\begin{align}
GP &= \text{Sales} - \text{Returns} - \text{Discounts} - \text{COGS} \\
\EBIT &= GP - (\text{Selling, General \& Administrative Expenses}) \\
\text{EBT} &= \EBIT - \text{Interest Expense} + \text{Non-Operating Income} \\
\text{Net Income} &= \text{EBT} - \text{Income Tax Expense}
\end{align}

\subsection{Vertical and Horizontal Analytical Models}
Forensic detection begins with relational normalizations:
\begin{enumerate}
    \item \textbf{Vertical Common-Size Transformation:} Converts absolute monetary values into scale-invariant percentages. A sudden expansion in Accounts Receivable from $12\%$ to $28\%$ of total assets signals aggressive revenue recognition or uncollectible receivables.
    \item \textbf{Horizontal Index Transformation:} Measures year-over-year rate of change:
    \begin{equation}
    g_{R} = \frac{\text{Sales}_t - \text{Sales}_{t-1}}{\text{Sales}_{t-1}}, \quad g_{AR} = \frac{\text{AR}_t - \text{AR}_{t-1}}{\text{AR}_{t-1}}
    \end{equation}
    Under healthy commercial dynamics, $g_{AR} \approx g_{R}$. When $g_{AR} \gg g_R$, it triggers the classic forensic red flag for channel stuffing or unearned revenue accruals.
\end{enumerate}

% ----------------------------------------------------------------------
% PART 3: VISUAL & GEOMETRIC INTUITION
% ----------------------------------------------------------------------
\section{Visual \& Geometric Intuition: The Fraud Diamond Architecture}

The behavioral and organizational prerequisites for corporate financial deceit are formalized through the \textbf{Fraud Triangle} (Donald Cressey) and its modern extension, the \textbf{Fraud Diamond} (David Wolfe and Dana Hermanson). Figure~\ref{fig:fraud_diamond} visualizes these four interacting vectors.

\begin{figure}[htbp]
\centering
\begin{tikzpicture}[
    scale=0.95,
    font=\sffamily\footnotesize,
    node distance=2.2cm,
    every node/.style={align=center}
]

% Diamond Nodes
\node[draw=coralborder, fill=coralbg, thick, rounded corners=4pt, text width=3.8cm, minimum height=1.3cm] (pressure) at (0, 3.2) {
    \textbf{1. INCENTIVE / PRESSURE}\\
    \textit{Earnings targets, stock option vesting, debt covenants, personal distress}
};

\node[draw=primarydark, fill=primarylight, thick, rounded corners=4pt, text width=3.8cm, minimum height=1.3cm] (opportunity) at (-4.2, 0) {
    \textbf{2. OPPORTUNITY}\\
    \textit{Weak internal controls, related-party deals, complex opaque accounting}
};

\node[draw=tealborder, fill=tealbg, thick, rounded corners=4pt, text width=3.8cm, minimum height=1.3cm] (rationalization) at (4.2, 0) {
    \textbf{3. RATIONALIZATION}\\
    \textit{``Just a temporary loan'', ``Everyone massages numbers'', ``We deserve it''}
};

\node[draw=purpleborder, fill=purplebg, thick, rounded corners=4pt, text width=3.8cm, minimum height=1.3cm] (capability) at (0, -3.2) {
    \textbf{4. CAPABILITY (DIAMOND)}\\
    \textit{Executive authority, technical prowess to coerce staff and deceive auditors}
};

% Central Target
\node[draw=red, fill=red!10, circle, thick, inner sep=6pt] (fraud) at (0, 0) {
    \textbf{\textcolor{red}{MATERIAL}}\\
    \textbf{\textcolor{red}{FINANCIAL}}\\
    \textbf{\textcolor{red}{FRAUD}}
};

% Arrows
\draw[<->, >=stealth, thick, color=coralborder] (pressure) -- (opportunity);
\draw[<->, >=stealth, thick, color=coralborder] (pressure) -- (rationalization);
\draw[<->, >=stealth, thick, color=purpleborder] (capability) -- (opportunity);
\draw[<->, >=stealth, thick, color=purpleborder] (capability) -- (rationalization);

\draw[->, >=stealth, line width=1.5pt, color=red] (pressure) -- (fraud);
\draw[->, >=stealth, line width=1.5pt, color=red] (opportunity) -- (fraud);
\draw[->, >=stealth, line width=1.5pt, color=red] (rationalization) -- (fraud);
\draw[->, >=stealth, line width=1.5pt, color=red] (capability) -- (fraud);

\end{tikzpicture}
\caption{The Fraud Diamond: Structural Synthesis of Pressure, Opportunity, Rationalization, and Executive Capability.}
\label{fig:fraud_diamond}
\end{figure}

% ----------------------------------------------------------------------
% PART 4: DEEP INTUITION & CONCEPTUAL DERIVATIONS
% ----------------------------------------------------------------------
\section{Deep Intuition \& Conceptual Derivations}

\subsection{The Mechanics of Earnings Manipulation Schemes}
Corporate financial statement manipulation generally follows two distinct strategies:

\subsubsection{1. Cookie-Jar Reserve Accounting}
During exceptionally profitable quarters, management artificially suppresses reported earnings by recognizing inflated estimated accruals (e.g., excessive provisions for bad debts, restructuring reserves, or inventory writedowns):
\begin{align*}
\text{Debit: } & \text{Bad Debt Expense} \quad \text{Rs.~10 Crores} \\
\text{Credit: } & \text{Allowance for Doubtful Accounts (Contra-Asset)} \quad \text{Rs.~10 Crores}
\end{align*}
In future lean quarters, when operating earnings slump, management reverses these reserves directly into income ($\text{Credit: Bad Debt Expense}$ or $\text{Operating Income}$), smoothing corporate performance and giving investors the illusion of low earnings volatility.

\subsubsection{2. Channel Stuffing and Bill-and-Hold Manipulations}
To satisfy aggressive quarterly revenue expectations, sales executives entice distributors to accept bulk orders beyond normal market absorption capacity, offering lenient return privileges, deferred payment terms, or deep price discounts:
\begin{align*}
\text{Debit: } & \text{Accounts Receivable} \quad \text{Rs.~50 Crores} \\
\text{Credit: } & \text{Revenue} \quad \text{Rs.~50 Crores}
\end{align*}
Because the customers hold right of return without economic risk, the transfer of control under IFRS 15 / ASC 606 is incomplete. In the subsequent quarter, massive inventory returns occur, triggering severe top-line collapse.

% ----------------------------------------------------------------------
% PART 5: EXPLICIT GOTCHAS & DIAGNOSTIC MISCONCEPTIONS
% ----------------------------------------------------------------------
\section{Explicit Gotchas \& Diagnostic Misconceptions}

\begin{trapbox}{Gotcha 1: Confusing Liquidity with Solvency}
\textbf{The Trap:} Assuming a firm with a high Current Ratio ($CA/CL > 3.0$) is immune to bankruptcy.\\
\textbf{Forensic Reality:} Current ratio can be heavily contaminated by illiquid, uncollectible, or fictitious current assets. An enterprise carrying Rs.~$100$ Crores in obsolete inventory and Rs.~$80$ Crores in uncollectible receivables against Rs.~$50$ Crores in current payables exhibits a Current Ratio of $3.6$, yet faces catastrophic insolvency when short-term debts mature.
\end{trapbox}

\begin{trapbox}{Gotcha 2: Treating EBITDA as a Proxy for Cash Flow}
\textbf{The Trap:} Evaluating corporate debt service capability using EBITDA ($\EBIT + \text{D\&A}$) under the belief that ``EBITDA equals cash earnings''.\\
\textbf{Forensic Reality:} EBITDA ignores working capital investments (e.g., cash tied up in soaring unpaid customer receivables), mandatory tax outflows, interest payments, and required replacement capital expenditures. As Warren Buffett famously noted: \textit{``Does management think the tooth fairy pays for capital expenditures?''}
\end{trapbox}

\begin{trapbox}{Gotcha 3: The Danger of Common-Size Trend Reversals}
\textbf{The Trap:} Relying exclusively on vertical common-size percentages when top-line revenues are themselves artificially inflated.\\
\textbf{Forensic Reality:} If an executive inflates revenue by $40\%$ through fictitious bill-and-hold invoices, an operating expense line that grew by $30\%$ in absolute terms will deceptively appear to have \textit{shrunk} as a percentage of revenue ($30\% / 140\% < \text{baseline}$). Vertical analysis must always be cross-calibrated against horizontal physical-unit indices.
\end{trapbox}

\begin{trapbox}{Gotcha 4: Rationalization vs.~Intent in Forensic Legal Proceedings}
\textbf{The Trap:} Assuming that proving executive rationalization (e.g., ``the CEO believed the market would recover next quarter'') exonerates the perpetrator of fraudulent misstatement.\\
\textbf{Forensic Reality:} In criminal jurisprudence, rationalization is a psychological enabler, not a legal defense. Intent (\textit{mens rea}) is established by the deliberate act of booking entries known to lack economic substance, regardless of whether the executive hoped to fix the shortfall later.
\end{trapbox}

% ----------------------------------------------------------------------
% PART 6: COMPARATIVE SYNTHESIS & SUMMARY TABLE
% ----------------------------------------------------------------------
\section{Comparative Synthesis \& Summary Table}

Table~\ref{tab:comparison_fraud_categories} compares the three core branches of the ACFE Fraud Tree across frequency, median financial loss, detection mechanisms, and algorithmic forensic signals.

\begin{table}[htbp]
\centering
\small
\begin{tabularx}{\textwidth}{lXXX}
\toprule
\textbf{Forensic Metric} & \textbf{Asset Misappropriation} & \textbf{Corruption Schemes} & \textbf{Financial Statement Fraud} \\
\midrule
\textbf{Incidence Rate} & Most frequent ($\approx 85\%$ of corporate fraud cases). & Intermediate frequency ($\approx 35\%$ of corporate fraud cases). & Least frequent ($\approx 10\%$ of corporate fraud cases). \\
\addlinespace
\textbf{Median Financial Loss} & Smallest loss ($\approx \$100{,}000$ median loss per scheme). & Moderate loss ($\approx \$250{,}000$ median loss per scheme). & Catastrophic loss ($\approx \$1{,}000{,}000+$ to enterprise collapse). \\
\addlinespace
\textbf{Primary Perpetrators} & Lower-level staff, payroll clerks, procurement officers. & Middle management, purchasing agents, contract negotiators. & C-Suite executives (CEO, CFO, Controllers, Board Promoters). \\
\addlinespace
\textbf{Core Mechanism} & Ghost employees, skimming cash, billing schemes, theft of inventory. & Bribery, kickbacks, conflicts of interest, illegal bid rigging. & Fictitious revenues, hidden liabilities, cookie jars, asset overvaluation. \\
\addlinespace
\textbf{Forensic Analytics} & Duplicate payments, ghost vendor phone/address matches. & Vendor concentration analysis, shell company beneficial ownership graph. & Horizontal/vertical statement variances, Beneish M-Score, accrual divergences. \\
\bottomrule
\end{tabularx}
\caption{Systemic Comparison of Corporate Fraud Categories (ACFE Taxonomy).}
\label{tab:comparison_fraud_categories}
\end{table}

% ----------------------------------------------------------------------
% PART 7: PYTHON / SCIENTIFIC COMPUTING LAB
% ----------------------------------------------------------------------
\section{Python / Scientific Computing Lab: Vertical \& Horizontal Statement Engine}

Listing~\ref{lst:financial_statement_analysis} provides an automated Python engine for processing multi-year financial statements, computing common-size distributions, evaluating horizontal expansion rates, and highlighting divergence red flags.

\begin{lstlisting}[language=Python, caption={Vertical and Horizontal Financial Statement Analysis in Python}, label={lst:financial_statement_analysis}]
import pandas as pd
import numpy as np

def analyze_financial_statements(fs_df: pd.DataFrame, revenue_col: str = 'Revenue') -> pd.DataFrame:
    """
    Performs forensic vertical (common-size) and horizontal (YoY growth) 
    analysis on multi-year financial statement items.
    
    Parameters:
        fs_df (pd.DataFrame): DataFrame indexed by line items, with fiscal year columns.
        revenue_col (str): Line item name representing total scale baseline.
        
    Returns:
        pd.DataFrame: Augmented DataFrame with vertical shares, growth rates, and flags.
    """
    years = [col for col in fs_df.columns if col.startswith('FY')]
    years.sort()
    
    results = fs_df.copy()
    
    # 1. Vertical Analysis (Common-Size % of Revenue)
    for yr in years:
        total_rev = fs_df.loc[revenue_col, yr]
        results[f'{yr}_Vertical%'] = (fs_df[yr] / total_rev * 100.0).round(2)
        
    # 2. Horizontal Analysis (Year-over-Year % Growth)
    for i in range(1, len(years)):
        prev_yr, curr_yr = years[i-1], years[i]
        results[f'YoY_Growth_{curr_yr}%'] = (
            (fs_df[curr_yr] - fs_df[prev_yr]) / fs_df[prev_yr] * 100.0
        ).round(2)
        
    # 3. Forensic Divergence Check: Flag AR growth exceeding Revenue growth by > 15%
    if 'Accounts_Receivable' in fs_df.index:
        for i in range(1, len(years)):
            prev_yr, curr_yr = years[i-1], years[i]
            rev_growth = results.loc[revenue_col, f'YoY_Growth_{curr_yr}%']
            ar_growth = results.loc['Accounts_Receivable', f'YoY_Growth_{curr_yr}%']
            
            flag_col = f'Flag_AR_Divergence_{curr_yr}'
            results[flag_col] = False
            if (ar_growth - rev_growth) > 15.0:
                results.loc['Accounts_Receivable', flag_col] = True
                
    return results

# -------------------------------------------------------------
# Test Dataset with Channel Stuffing Anomaly
# -------------------------------------------------------------
if __name__ == '__main__':
    statement_data = {
        'FY2023': [1000.0, 600.0, 400.0, 150.0, 250.0, 120.0, 50.0],
        'FY2024': [1100.0, 660.0, 440.0, 165.0, 275.0, 140.0, 55.0],
        'FY2025': [1250.0, 720.0, 530.0, 180.0, 350.0, 260.0, 45.0]  # AR spikes to 260!
    }
    line_items = [
        'Revenue', 'COGS', 'Gross_Profit', 'SG&A', 'EBIT', 
        'Accounts_Receivable', 'Operating_Cash_Flow'
    ]
    df = pd.DataFrame(statement_data, index=line_items)
    audit_report = analyze_financial_statements(df)
    
    print("Forensic Vertical and Horizontal Analysis:")
    print(audit_report[['FY2024_Vertical%', 'FY2025_Vertical%', 'YoY_Growth_FY2025%']])
    if 'Flag_AR_Divergence_FY2025' in audit_report.columns:
        print("\nAR Divergence Flagged in FY2025:", 
              audit_report.loc['Accounts_Receivable', 'Flag_AR_Divergence_FY2025'])
\end{lstlisting}

% ----------------------------------------------------------------------
% PART 8: COMPREHENSIVE PRACTICE PROBLEMS
% ----------------------------------------------------------------------
\section{Comprehensive Practice Problems}
\textit{Complete mathematical derivations and diagnostic MCQ distractor rationales are located in the Appendix.}

\subsection*{Part A: Analytical \& Computational Problems}

\begin{enumerate}[label=\textbf{P2.\arabic*}]
    \item \textbf{Multi-Step Income Statement Reconciliation:}\\
    An Indian telecommunications enterprise presents the following trial balance extracts for the financial year (all figures in Rs.~Crores):
    \begin{center}
    \begin{tabular}{lr|lr}
    \toprule
    \textbf{Account} & \textbf{Debit} & \textbf{Account} & \textbf{Credit} \\
    \midrule
    Gross Sales Revenue & --- & Gross Sales Revenue & $2{,}400$ \\
    Sales Returns \& Allowances & $80$ & Non-Operating Interest Income & $35$ \\
    Cost of Goods Sold (COGS) & $1{,}240$ & Gain on Disposal of Tower Assets & $25$ \\
    Administrative Expenses & $220$ & Equity Capital & $500$ \\
    Sales \& Marketing Expenses & $180$ & Long-Term Debt ($9\%$ coupon) & $400$ \\
    Depreciation on Capital Assets & $140$ & Trade Payables & $190$ \\
    Corporate Income Tax Expense ($25\%$) & ? & Retained Earnings & $750$ \\
    \bottomrule
    \end{tabular}
    \end{center}
    \begin{enumerate}[label=(\alph*)]
        \item Compute Net Sales Revenue, Gross Profit, and Operating Profit (EBIT).
        \item Compute EBITDA and Net Profit After Tax.
        \item If management capitalized Rs.~$60$ Crores of marketing costs into intangible assets during the year, reconstruct true EBIT and quantify the percentage overstatement.
    \end{enumerate}

    \item \textbf{Working Capital \& Window Dressing Forensics:}\\
    At fiscal year-end (March 30), a consumer electronics firm reports Current Assets of Rs.~$60{,}00{,}000$ and Current Liabilities of Rs.~$30{,}00{,}000$ (Current Ratio $= 2.0$). To satisfy a bank loan covenant requiring a Current Ratio $\ge 2.2$, the CFO borrows Rs.~$20{,}00{,}000$ in short-term bank overdraft on March 31 and uses Rs.~$15{,}00{,}000$ of available cash to settle trade payables.
    \begin{enumerate}[label=(\alph*)]
        \item Compute the final Current Ratio on March 31 following these actions.
        \item Did the CFO's actions successfully meet the covenant threshold?
        \item Explain why paying off liabilities with cash mathematically improves a Current Ratio that is initially greater than $1.0$.
    \end{enumerate}

    \item \textbf{Common-Size Horizontal Forensic Analysis:}\\
    Over a three-year period, a pharmaceutical company reports the following financial figures (in Rs.~Lakhs):
    \begin{center}
    \begin{tabular}{lrrr}
    \toprule
    \textbf{Metric} & \textbf{Year 1} & \textbf{Year 2} & \textbf{Year 3} \\
    \midrule
    Sales Revenue & $5{,}000$ & $5{,}800$ & $6{,}400$ \\
    Accounts Receivable & $600$ & $950$ & $1{,}600$ \\
    Inventory & $800$ & $1{,}100$ & $1{,}750$ \\
    Allowance for Bad Debts & $60$ & $40$ & $30$ \\
    \bottomrule
    \end{tabular}
    \end{center}
    \begin{enumerate}[label=(\alph*)]
        \item Compute the horizontal percentage growth of Sales, Accounts Receivable, and Inventory from Year 1 to Year 3.
        \item Calculate the Days Sales Outstanding (DSO) for each year assuming a 365-day year.
        \item Analyze the forensic red flags revealed by the trajectory of the Allowance for Bad Debts relative to Receivables.
    \end{enumerate}

    \item \textbf{The Fraud Diamond in Corporate Governance:}\\
    Evaluate the Enron collapse through the lens of the Fraud Diamond. Provide concrete historical evidence illustrating how Incentive, Opportunity, Rationalization, and Capability manifested among Enron's senior leadership (Ken Lay, Jeff Skilling, and Andy Fastow).

    \item \textbf{Off-Balance Sheet Debt \& Operating Leases:}\\
    Prior to IFRS 16 / ASC 842, an airline entered into long-term aircraft operating leases requiring annual lease payments of Rs.~$100$ Crores for the next 10 years. The company's incremental borrowing cost is $8\%$.
    \begin{enumerate}[label=(\alph*)]
        \item Calculate the unrecorded present value of off-balance sheet lease commitments.
        \item Explain how keeping these obligations off the balance sheet distorted the firm's Debt-to-Equity and Return-on-Assets ratios.
    \end{enumerate}
\end{enumerate}

\subsection*{Part B: Diagnostic Multiple-Choice Questions}

\begin{enumerate}[label=\textbf{MCQ 2.\arabic*}]
    \item A manufacturing enterprise artificially inflates its reported Gross Profit margin during Year 2. Forensic investigation reveals that physical production volumes remained constant, unit selling prices did not change, and raw material costs were steady. Which of the following accounting maneuvers directly produces this fraudulent outcome?
    \begin{enumerate}[label=(\alph*)]
        \item Improperly shifting closing inventory valuation from FIFO to a fictitious higher ending inventory balance.
        \item Reclassifying administrative office utility expenses into factory selling expenses.
        \item Writing off obsolete spare parts as an extraordinary impairment loss in the income statement.
        \item Accelerating depreciation on executive office headquarters from straight-line to double-declining balance.
    \end{enumerate}

    \item In the context of the Fraud Diamond, which factor distinguishes Wolfe and Hermanson's model from Cressey's classical Fraud Triangle?
    \begin{enumerate}[label=(\alph*)]
        \item The presence of external financial pressure, such as impending debt covenant defaults.
        \item Executive capability, encompassing technical skill, interpersonal influence, and authoritarian access required to execute and conceal deceit.
        \item Moral rationalization by the perpetrator that the stolen funds are merely a short-term loan.
        \item Deficiencies in corporate governance and an absence of segregation of duties.
    \end{enumerate}
\end{enumerate}
